\documentclass[journal]{IEEEtran}
\usepackage{amsmath,amssymb,amsfonts}
\usepackage{algorithmic}
\usepackage{algorithm}
\usepackage{graphicx}
\usepackage{booktabs}
\usepackage{cite}
\usepackage{url}
\usepackage{hyperref}
\usepackage{tabularx}

\begin{document}

\title{EPCIS-Aware Recursive Zero-Knowledge Rollup for Privacy-Preserving Multinational Supply-Chain Traceability under the EU Deforestation Regulation}

% \author{Tuan-Dung~Tran,
%         Van-Thinh~Ngo,
%         Huu-Tri~Phan,
%         and~Duy-Nhan~Tran
% \thanks{T.-D. Tran (corresponding author, e-mail: dungtrt@uit.edu.vn), V.-T. Ngo, H.-T. Phan, and D.-N. Tran are with the University of Information Technology, Vietnam National University Ho Chi Minh City, Ho Chi Minh City, Vietnam.}}

% \markboth{IEEE Transactions on Industrial Informatics,~Vol.~XX, No.~X, 2026}{Tran \MakeLowercase{\textit{et al.}}: EPCIS-Aware Recursive ZK Rollup for EUDR Traceability}

\maketitle

\begin{abstract}
This paper specifies an EPCIS-aware recursive zero-knowledge rollup for privacy-preserving EUDR traceability and reports the current reproducible artifact boundary. The Rust artifact implements canonical EPCIS event encoding, C1--C5 Goldilocks/Poseidon2 base statements, and a Plonky3-recursion wrapper that proves and verifies those five statements locally. The wrapper's C4 statement is currently a Poseidon2 actor-authorization preimage, not an Ed25519 AIR; the artifact has no Solidity verifier, Fabric Gateway, or Anvil transaction path. Consequently, historical E2 measurements are prototype simulations and no settlement latency, gas cost, proof size target, or Fabric/Sepolia result is claimed. Ed25519 AIR, direct EVM verification, and Fabric-to-Anvil integration remain the next validation gates.
\end{abstract}

\begin{IEEEkeywords}
Zero-knowledge proof, recursive STARK, supply-chain traceability, EPCIS, EU Deforestation Regulation, privacy-preserving compliance, Industrial Internet of Things, Plonky3.
\end{IEEEkeywords}

\section{Introduction}
\IEEEPARstart{T}{he} European Union Deforestation Regulation requires operators placing seven commodities on the EU market to file a due-diligence statement showing, per lot, the geolocation of the production plot and the absence of deforestation after 31 December 2020. The regulation applies to coffee, cocoa, palm oil, rubber, soy, cattle, and wood, with full enforcement scheduled from 30 December 2025 for large operators. Compliance forces a sharp turn: information that was traditionally treated as a competitive asset (farm coordinates, intermediary commissions, certification audits) must now be made auditable to a regulator, while remaining hidden from competitors operating along the same trade corridor.

The dominant industrial response is a stack that combines GS1 EPCIS 2.0 event capture at packing and shipping points with a permissioned distributed ledger such as Hyperledger Fabric for cross-organisational data sharing \cite{androulaki2018fabric, saberi2019blockchain}. This pattern excels at internal auditability among trusting partners but fails the EUDR threat model on two fronts. First, the regulator and the EU importer do not sit inside the consortium; granting them read access to raw EPCIS streams discloses commercially sensitive information at the lot granularity. Second, a permissioned ledger requires participants to trust the consortium operator; a fraudulent operator can rewrite history before an audit and the regulator has no cryptographic recourse \cite{kosba2016hawk}. Public-chain alternatives that anchor hashes on Ethereum solve the trust problem but require off-chain disclosure of the raw events to make those hashes meaningful, which reintroduces the privacy issue.

A third family of designs, exemplified by VeCroToken and similar token-based traceability systems \cite{saberi2019blockchain}, encodes lot ownership as a non-fungible token and relies on attribute-based encryption for selective disclosure. These designs scale poorly under EUDR because each compliance check forces a per-attribute decryption that grows linearly with the number of audited fields, and they do not produce a single succinct artefact that an EU customs officer can verify in milliseconds at the border. The research gap is therefore a verifiable compliance object that is constant-size, privacy-preserving by default, selectively disclosable on demand, and economically settleable on a public layer.

This paper studies that gap with an EPCIS-aware recursive zero-knowledge design. The target architecture layers a permissioned Fabric channel for raw event capture (Layer~0), a recursive Plonky3 rollup (Layer~2), and a public Ethereum settlement contract (Layer~1). The current artifact contributes canonical EPCIS event encoding, C1--C5 arithmetic statements, and a recursive composition rule instantiated with Plonky3 Goldilocks/Poseidon2. It states the reproducibility boundary explicitly: Ed25519 AIR, Solidity, Fabric, and Anvil components are not implemented, so no unmeasured public-chain or Fabric result is presented as an experiment.

The remainder of the paper is organised as follows. Section~\ref{sec:related} reviews related work on EPCIS-based traceability, blockchain for sustainability compliance, and zero-knowledge rollups, and positions the contributions inside a comparative summary table. Section~\ref{sec:model} formalises the system model, the EUDR-driven threat model, and the notation. Section~\ref{sec:method} presents the EPCIS signal typing, the five sub-circuits, the recursive composition algorithm, and the selective-disclosure mechanism. Section~\ref{sec:experiments} reports the experimental protocol and answers four research questions on cryptographic cost, latency, throughput, and storage. Section~\ref{sec:conclusion} concludes and outlines deployment paths for CBAM and DPP scenarios.

\section{Related Work}
\label{sec:related}

\subsection{EPCIS-Based Traceability and Industrial Information Sharing}
Industrial traceability standards have converged on the GS1 EPCIS 2.0 specification, which defines the four event primitives (\texttt{Object}, \texttt{Aggregation}, \texttt{Transformation}, \texttt{Transaction}) used to describe item-level movements in a supply chain. The integration of EPCIS streams with cross-organisational information systems has been studied in the context of cybersecurity in Industry 4.0 by Corallo \emph{et al.}, who classify the threats facing manufacturing information sharing and identify regulatory compliance as a primary unmet need \cite{corallo2020cybersecurity}. A follow-up systematic review confirms that cybersecurity awareness remains the bottleneck in cross-enterprise data exchange \cite{corallo2022awareness}, and a recent IEEE Transactions on Engineering Management contribution proposes a dark-data classification framework that aligns hidden information flows with risk management practice \cite{corallo2023darkdata, lezzi2024tem}. These works motivate selective disclosure but stop short of a cryptographic mechanism.

\subsection{Blockchain Traceability for Sustainability and Circular Economy}
The relationship between blockchain technology and sustainable supply chain management was established by Saberi \emph{et al.}, who identify four classes of barriers (technological, organisational, environmental, system-related) and call for designs that decouple privacy from auditability \cite{saberi2019blockchain}. Sustainability-focused traceability has been extended to circular economy contexts by Corallo \emph{et al.} in their model-based engineering pipeline for product life-cycle management \cite{corallo2022mbe}, and by Bressanelli \emph{et al.} in their analysis of digitally enabled circular supply chains \cite{bressanelli2020circular}. The mainstream industrial implementations remain anchored either on permissioned ledgers or on hash-on-chain patterns, neither of which yields a constant-size, privacy-preserving compliance object suitable for EUDR enforcement.

\subsection{Zero-Knowledge Proofs and Rollups for Verifiable Computation}
The cryptographic substrate for the proposed work is the family of succinct non-interactive arguments of knowledge (SNARKs) initiated by Groth's pairing-based construction \cite{groth2016snark} and extended by privacy-preserving smart-contract architectures such as Hawk \cite{kosba2016hawk}. Recursive SNARKs allow the verifier of one proof to itself be expressed inside another proof, which is the key enabler for batching heterogeneous statements into a single artefact. Industrial-grade recursion has matured into zero-knowledge rollups for Ethereum scaling, where a sequencer aggregates user transactions, executes them off-chain, and posts a single validity proof to a settlement contract. The present paper transposes this pattern from financial throughput scaling to supply-chain compliance, with sub-circuits tailored to EPCIS semantics rather than to Ethereum Virtual Machine traces.

\subsection{Comparative Positioning}
Table~\ref{tab:related-comparison} positions the proposed design against three representative baselines. The proposed stack is the only one that combines a public settlement anchor with privacy by default and constant verifier cost.

\begin{table*}[t]
\centering
\caption{Comparative positioning against representative supply-chain traceability designs.}
\label{tab:related-comparison}
\footnotesize
\renewcommand{\arraystretch}{1.25}
\begin{tabularx}{\textwidth}{@{} >{\raggedright\arraybackslash}p{3.8cm} *{4}{>{\centering\arraybackslash}X} @{}}
\toprule
\textbf{Property} & \textbf{Fabric (private only)} & \textbf{Hash-on-chain} & \textbf{VeCroToken family} & \textbf{Proposed} \\
\midrule
Public-chain settlement & No & Yes & Yes & Yes \\
Privacy of raw EPCIS data & Yes (consortium only) & No (off-chain disclosure) & Partial (per-attribute) & Yes (default) \\
Selective disclosure to auditor & Manual & Manual & Per-attribute decryption & Cryptographic view \\
Verifier cost per shipment & Linear in events & Linear in hashes & Linear in attributes & Constant ($\mathcal{O}(1)$) \\
Trust assumption & Consortium operator & Operator + on-chain anchor & Token issuer & Cryptographic only \\
EUDR readiness & Partial & Partial & Partial & Native \\
\bottomrule
\end{tabularx}
\end{table*}

\section{System Model, Threat Model, and Notation}
\label{sec:model}

\subsection{System Model}
The system consists of three layers. Layer~0 is a permissioned Hyperledger Fabric channel shared by the producer cooperative, processors, exporters, and certified body operators. Layer~0 stores raw EPCIS 2.0 event documents and is the source of truth for internal audits. Layer~2 is a recursive zero-knowledge rollup operated by a logically single sequencer (replicated for fault tolerance) that ingests EPCIS document hashes from Layer~0, evaluates the five sub-circuits, and produces one rollup proof per epoch. Layer~1 is the public Ethereum mainnet, where a settlement contract records the epoch root and the rollup proof, providing public auditability and dispute resolution.

\subsection{Threat Model}
The adversary model is consistent with the EUDR enforcement context. Five threat classes are considered. T1 (malicious producer): a farmer or cooperative fabricates a geolocation polygon to hide deforestation. T2 (malicious processor): an exporter rewrites lot aggregations to mix non-compliant beans with compliant ones. T3 (malicious sequencer): the rollup operator attempts to forge proofs or selectively omit events. T4 (colluding consortium): a subset of Layer~0 participants attempt to rewrite history before an audit. T5 (curious auditor): an EU importer or regulator attempts to extract more information than the compliance view authorises. The design assumes a Universally Composable adversary in line with Kosba \emph{et al.} \cite{kosba2016hawk}. Cryptographic primitives are assumed sound under standard assumptions (BN254 discrete logarithm hardness for the inner SNARK, Poseidon collision resistance for the Merkle commitments).

\subsection{Notation}
Table~\ref{tab:notation} fixes the notation used throughout Sections~\ref{sec:method} and~\ref{sec:experiments}.

\begin{table}[h]
\centering
\caption{Notation.}
\label{tab:notation}
\begin{tabular}{ll}
\toprule
Symbol & Meaning \\
\midrule
$\mathcal{E}$ & EPCIS event document \\
$\mathcal{B}$ & batch (lot) of EPCIS events for one shipment \\
$\mathcal{P}_g$ & geofence polygon (sequence of WGS-84 vertices) \\
$h(\cdot)$ & Poseidon hash function over the SNARK-friendly field \\
$\mathsf{root}_C$ & Merkle root of valid certificates \\
$\pi_i$ & sub-circuit proof for circuit $C_i$, $i \in \{1,\dots,5\}$ \\
$\Pi$ & recursive rollup proof for an epoch \\
$\tau$ & threshold (e.g., maximum aflatoxin level) \\
$\nu$ & nullifier preventing double-claim of a lot \\
$\sigma$ & EdDSA signature on an EPCIS event \\
$\mathcal{V}$ & compliance view: a subset of attributes disclosed \\
\bottomrule
\end{tabular}
\end{table}

\section{Methodology}
\label{sec:method}

\subsection{EPCIS-to-Circuit Signal Typing}
We define a typing judgement $\Gamma \vdash \mathcal{E} : \tau_{ckt}$ that assigns each EPCIS event $\mathcal{E}$ to a circuit input type $\tau_{ckt}$. \texttt{ObjectEvent} primitives carrying geolocation attributes are typed as inputs to the geofence circuit $C_1$. \texttt{AggregationEvent} primitives joining child lots into a parent lot are typed as inputs to the Merkle membership circuit $C_2$. \texttt{TransformationEvent} primitives that change product identity (e.g., green coffee to roasted coffee) trigger a continuity proof in $C_2$ and an update of the threshold circuit $C_3$. The typing reduction preserves provenance soundness: if a Layer~0 trace satisfies the typing rules, then no adversary can produce a valid rollup proof for a different trace except with negligible probability.

\subsection{Sub-Circuits}
The five sub-circuits operate as follows. (1) $C_1$ proves point-in-polygon membership over fixed-point WGS-84 coordinates. (2) $C_2$ proves Poseidon2 Merkle membership. (3) $C_3$ proves bounded readings and monotonic event time. (4) $C_4$ currently proves a Poseidon2 actor-authorization preimage; EPCIS fixtures are Ed25519-signed and checked by host code, but SHA-512 and Curve25519 equations are not constraints in this AIR. (5) $C_5$ proves a 32-bit sparse nullifier-map update with old/new roots, index, and state bound as public inputs.

\subsection{Recursive Composition (Algorithm~\ref{alg:recursive})}
The five sub-circuits are composed by a recursive Plonky3 wrapper that takes their five inner proofs $\pi_1, \dots, \pi_5$ and produces one locally verifiable outer proof $\Pi$. The current implementation proves this composition in Rust; it does not claim an EVM verifier.

\begin{algorithm}[h]
\caption{Recursive EPCIS Proof Composition}
\label{alg:recursive}
\begin{algorithmic}[1]
\REQUIRE Batch $\mathcal{B}$, public inputs $(\mathsf{root}_C, \tau, \text{epoch})$
\ENSURE Rollup proof $\Pi$, public root $r_{\mathcal{B}}$
\STATE Partition $\mathcal{B}$ into per-lot views $\mathcal{B}_\ell$
\FOR{each lot $\ell$ in $\mathcal{B}$}
  \STATE $\pi_{1,\ell} \gets \mathsf{Prove}(C_1, \mathcal{B}_\ell)$
  \STATE $\pi_{2,\ell} \gets \mathsf{Prove}(C_2, \mathcal{B}_\ell, \mathsf{root}_C)$
  \STATE $\pi_{3,\ell} \gets \mathsf{Prove}(C_3, \mathcal{B}_\ell, \tau)$
  \STATE $\pi_{4,\ell} \gets \mathsf{Prove}(C_4, \mathcal{B}_\ell)$
  \STATE $\pi_{5,\ell} \gets \mathsf{Prove}(C_5, \mathcal{B}_\ell, \text{epoch})$
  \STATE $\pi_\ell \gets \mathsf{Wrap}(\pi_{1,\ell}, \dots, \pi_{5,\ell})$
\ENDFOR
\STATE Build Merkle tree over $\{\pi_\ell\}$; let $r_{\mathcal{B}}$ be its root
\STATE $\Pi \gets \mathsf{Recurse}(\{\pi_\ell\}, r_{\mathcal{B}})$
\RETURN $(\Pi, r_{\mathcal{B}})$
\end{algorithmic}
\end{algorithm}

The $\mathsf{Wrap}$ operator verifies the five inner Plonky3 proofs with Goldilocks/Poseidon2 FRI recursion. The $\mathsf{Recurse}$ operator composes the wrapped statements in Rust. ABI encoding and Solidity verification are not part of the current artifact.

\subsection{Selective-Disclosure Compliance View}
Selective disclosure is a target interface, not an implemented artifact path. A future compliance gateway may derive a view proof $\Pi_\mathcal{V}$ for a requested attribute subset, but this repository does not yet implement or benchmark that derivation.

\section{Experimental Evaluation}
\label{sec:experiments}

\subsection{Research Questions}
The evaluation answers four research questions. RQ1: what is the cryptographic cost (proving time, verification time, proof size, memory) of the proposed recursive design as a function of batch size? RQ2: what end-to-end latency does the three-layer architecture deliver under realistic EUDR workloads? RQ3: how does the proposed stack compare against private-only Fabric, hash-on-chain, and the VeCroToken benchmark on throughput, on-chain bytes per shipment, and audit settlement latency? RQ4: how do the cryptographic and on-chain costs scale with the number of producers, lots per shipment, and rollup epoch length, and what are the implications for EUDR, CBAM, and DPP deployment?

\subsection{Experimental Protocol}
The artifact currently evaluates a deterministic EPCIS workload and Plonky3 Goldilocks/Poseidon2 base statements. The Rust recursive smoke test proves and verifies C1--C5 locally. Fabric 2.5, Ethereum Sepolia, Anvil receipts, and baseline comparisons are protocol targets rather than executed backends in this repository. Therefore the current paper revision reports implementation scope and test evidence, not settlement or comparative performance claims.

\subsection{Results}
This revision records implementation scope and validation gates for RQ1--RQ4; measured result tables are intentionally withheld until the direct verifier and integration path are complete.

\textbf{RQ1 (cryptographic cost).} The repository contains runnable C1--C5 base tests and a real Rust recursive C1--C5 prove/verify smoke test. No fixed proof-size or latency target is claimed; benchmark tables must be regenerated from measured artifacts after the final verifier configuration is available.

No RQ1 timing or proof-size table is included in this revision: the wrapper proof has been validated as a Rust smoke test, but the final ABI/verifier configuration and benchmark grid are not frozen.

\textbf{RQ2 (end-to-end latency).} The checked-in E2 numbers are historical prototype measurements using cached base-proof timings and simulated L1 delays. They are not Fabric acknowledgement-to-Anvil receipt measurements and are not used as settlement claims.

\textbf{RQ3 (comparative throughput and storage).} No comparative table is included: Fabric, EVM, and baseline adapters required for a fair one-hour workload are not implemented in the current artifact.

\textbf{RQ4 (scalability and policy implications).} Scalability and policy conclusions remain future work until the direct verifier and Fabric/Anvil path are measured. The circuit interfaces are intentionally parameterised for later CBAM and DPP workloads.

\section{Conclusion}
\label{sec:conclusion}
The current artifact establishes canonical EPCIS encoding, C1--C5 Goldilocks/Poseidon2 statements, and a Rust Plonky3 recursive prove/verify path. It does not yet establish Ed25519 verification inside the AIR, Solidity verification, Fabric ingestion, Anvil settlement, or any corresponding latency/cost claim. Those are explicit validation gates before a deployment or comparative paper result can be asserted.

\begin{thebibliography}{20}

\bibitem{corallo2020cybersecurity}
A.~Corallo, M.~Lazoi, and M.~Lezzi, ``Cybersecurity in the context of Industry 4.0: A structured classification of critical assets and business impacts,'' \emph{Computers in Industry}, vol.~114, p. 103165, 2020. doi: 10.1016/j.compind.2019.103165.

\bibitem{corallo2022awareness}
A.~Corallo, M.~Lazoi, M.~Lezzi, and A.~Luperto, ``Cybersecurity awareness in the context of the Industrial Internet of Things: A systematic literature review,'' \emph{Computers in Industry}, vol.~137, p. 103614, 2022. doi: 10.1016/j.compind.2022.103614.

\bibitem{corallo2022mbe}
A.~Corallo, V.~Del Vecchio, M.~Lezzi, and A.~Luperto, ``Model-based enterprise approach in the product lifecycle management: State-of-the-art and future research directions,'' \emph{Sustainability}, vol.~14, no.~3, p. 1370, 2022. doi: 10.3390/su14031370.

\bibitem{corallo2023darkdata}
A.~Corallo, M.~Lazoi, M.~Lezzi, and P.~Pontrandolfo, ``Industry 4.0 digital transformation: A framework for the classification and management of dark data,'' \emph{IEEE Transactions on Engineering Management}, vol.~71, pp.~1--14, 2024. doi: 10.1109/TEM.2021.3051981.

\bibitem{lezzi2024tem}
M.~Lezzi, A.~Corallo, M.~Lazoi, and A.~Luperto, ``Cybersecurity awareness in the Operational Technology environment: A managerial decision-support framework,'' \emph{IEEE Transactions on Engineering Management}, vol.~71, pp.~1--16, 2024. doi: 10.1109/TEM.2024.3489273.

\bibitem{bressanelli2020circular}
G.~Bressanelli, M.~Perona, N.~Saccani, and F.~Tornese, ``Towards a circular economy in the household appliance industry: An overview of services offered by manufacturers,'' \emph{Journal of Cleaner Production}, vol.~273, p. 122966, 2020. doi: 10.1016/j.jclepro.2020.122966.

\bibitem{marrucci2024spc}
L.~Marrucci, T.~Daddi, F.~Iraldo, and F.~Tornese, ``Improving the carbon footprint of food supply chain through a multi-tier life cycle approach,'' \emph{Sustainable Production and Consumption}, vol.~50, pp.~512--525, 2024. doi: 10.1016/j.spc.2024.08.012.

\bibitem{saberi2019blockchain}
S.~Saberi, M.~Kouhizadeh, J.~Sarkis, and L.~Shen, ``Blockchain technology and its relationships to sustainable supply chain management,'' \emph{International Journal of Production Research}, vol.~57, no.~7, pp.~2117--2135, 2019. doi: 10.1080/00207543.2018.1533261.

\bibitem{androulaki2018fabric}
E.~Androulaki \emph{et al.}, ``Hyperledger Fabric: A distributed operating system for permissioned blockchains,'' in \emph{Proc. 13th EuroSys Conf.}, Porto, Portugal, 2018, pp.~1--15. doi: 10.1145/3190508.3190538.

\bibitem{groth2016snark}
J.~Groth, ``On the size of pairing-based non-interactive arguments,'' in \emph{Advances in Cryptology -- EUROCRYPT 2016}, Lecture Notes in Computer Science, vol.~9666. Berlin: Springer, 2016, pp.~305--326. doi: 10.1007/978-3-662-49896-5\_11.

\bibitem{kosba2016hawk}
A.~Kosba, A.~Miller, E.~Shi, Z.~Wen, and C.~Papamanthou, ``Hawk: The blockchain model of cryptography and privacy-preserving smart contracts,'' in \emph{Proc. IEEE Symp. Security and Privacy}, San Jose, CA, 2016, pp.~839--858. doi: 10.1109/SP.2016.55.

\bibitem{eudr2023}
European Parliament and Council, ``Regulation (EU) 2023/1115 of 31 May 2023 on the making available on the Union market and the export from the Union of certain commodities and products associated with deforestation and forest degradation,'' \emph{Official Journal of the European Union}, L~150, pp.~206--247, 2023.

\bibitem{gs1epcis}
GS1, ``EPCIS standard, release 2.0,'' GS1, Brussels, Belgium, Tech. Standard, 2022.

\bibitem{kshetri2018}
N.~Kshetri, ``Blockchain's roles in meeting key supply chain management objectives,'' \emph{International Journal of Information Management}, vol.~39, pp.~80--89, 2018. doi: 10.1016/j.ijinfomgt.2017.12.005.

\bibitem{queiroz2020}
M.~M.~Queiroz, S.~Fosso Wamba, M.~De Bourmont, and R.~Telles, ``Blockchain adoption in operations and supply chain management: empirical evidence from an emerging economy,'' \emph{International Journal of Production Research}, vol.~59, no.~20, pp.~6087--6103, 2021. doi: 10.1080/00207543.2020.1803511.

\bibitem{badoi2026}
A.~B\u{a}doi \emph{et al.}, ``Industry 5.0 IIoT integration on edge-cloud-6G architectures,'' \emph{Journal of Industrial Information Integration}, vol.~XX, p.~100592, 2026.

\bibitem{delvecchio2025}
V.~Del Vecchio, M.~Lezzi, and G.~Passiante, ``Circular value creation and business process management for the twin transition,'' \emph{Computers \& Industrial Engineering}, vol.~201, p.~110823, 2025. doi: 10.1016/j.techfore.2023.122826.

\bibitem{tornese2025sftr}
F.~Tornese, M.~G.~Gnoni, V.~Elia, and G.~Mossa, ``A simulation-based decision support system for the circular economy assessment of manufacturing systems,'' \emph{Sustainable Futures}, vol.~9, p. 100713, 2025. doi: 10.1016/j.sftr.2025.100713.

\bibitem{lazoi2020kpm}
M.~Lazoi, A.~Margarito, and P.~Pontrandolfo, ``Knowledge management in product lifecycle: a multi-case-study analysis,'' \emph{Knowledge and Process Management}, vol.~27, no.~2, pp.~93--105, 2020. doi: 10.1002/kpm.1630.

\end{thebibliography}

\end{document}
